\section{Motivacija}
Kriptografija obezbedjuje poverljivost, integritet i autenticnost podataka.
U ovom radu koristimo standardni zapis i okruzenja za matematicke izraze.

\subsection{Osnovni pojmovi}
\begin{definition}
Kriptosistem je uredjeni skup $(P, C, K, E, D)$ gde su $P$ skup otvorenih
tekstova, $C$ skup sifrovanih tekstova, $K$ skup kljuceva, a $E$ i $D$
funkcije enkripcije i dekripcije. 
\end{definition}

\begin{theorem}
Za svako $m \in P$ i $k \in K$ vazi $D_k(E_k(m)) = m$.
\end{theorem}

\begin{proof}
Sledi direktno iz definicije korektnosti kriptosistema.
\end{proof}

\begin{example}
Za Cezarovu sifru sa pomerajem $k=3$, slovo $A$ prelazi u slovo $D$.
\end{example}

\section{Ilustracija}
Slika \ref{fig:primer} prikazuje mesto za ilustraciju. Ako fajl ne postoji,
slika se preskace bez prekida kompilacije.

\begin{figure}[H]
	\centering
	\IfFileExists{images/primer.pdf}{%
		\includegraphics[width=0.7\linewidth]{primer.pdf}%
	}{%
		\IfFileExists{images/primer.png}{%
			\includegraphics[width=0.7\linewidth]{primer.png}%
		}{}
	}
	\caption{Primer slike}
	\label{fig:primer}
\end{figure}


Knjiga \cite{smart2016cryptography} je dobar uvod u kriptografiju.